% ============================================================================================================
\chapter{Revue critique}\label{ch:revue}
% ============================================================================================================

Cette revue relit le travail de la phase 3.1 avec les règles fixées au départ : ne pas modifier les paramètres pour
améliorer les résultats, ne pas conserver une hypothèse parce qu'elle est dans le code, ne pas confondre calibration et
validation, ne pas masquer les limites des données, ne pas chercher des résultats séduisants, ne pas conserver les
conclusions de la v0.7 si elles sont invalidées.

\section{Ce qui a été fait conformément au mandat}

\begin{itemize}
\item \textbf{Paramètre de transmission.} ρ et θ ont été choisis par un critère externe (durée des écarts de production
observée, ancrage sur la consommation), avant de regarder les stratégies. Le choix dégrade nettement les résultats
budgétaires des stratégies d'offre (S2, S3, S5) : il ne les rend pas plus séduisants. Une plage est publiée, faute
d'estimation ponctuelle robuste.
\item \textbf{Calibration et validation.} Le backtest réestime les équations sur les seules données disponibles à
chaque date d'origine et utilise un potentiel construit en temps réel ; les résultats défavorables (dette surestimée de
20 points en 2019, chômage mal suivi, modèle battu par la prévision naïve sur 2010-2019) sont publiés.
\item \textbf{Conclusions de la v0.7 révisées.} Quatre conclusions de la v0.7 sont explicitement invalidées
(chapitre~\ref{ch:strategies}).
\item \textbf{Défauts trouvés et corrigés en cours de route} : le conflit de noms de la v0.7 (amortissement), deux
comparaisons inéquitables (une variante de ρ comparée à la référence centrale ; des variantes de backtest avec des
politiques différentes), et, dans la distribution dynamique, trois défauts de représentation des gagnants et perdants :
le relèvement de l'âge légal apparaissait comme un gain pour le premier décile ; la baisse des allocations chômage,
les transferts aux ménages et le point d'indice n'avaient pas de perdants ou de gagnants directs ; le contrôle de
cohérence comparait un revenu nominal macroéconomique (hors TVA) à une somme microéconomique (avec TVA). Toutes les
simulations concernées ont été relancées ; les résultats antérieurs sont conservés pour traçabilité.
\end{itemize}

\section{Ce qui reste fragile}

\begin{enumerate}
\item \textbf{Le backtest ne valide pas le modèle pour le long terme.} Sur 2010-2019, la dette est surestimée de 20
points en neuf ans, principalement parce que la productivité en temps réel était surestimée (erreur que commettaient
aussi les prévisionnistes de l'époque) et que la charge d'intérêts réagit trop lentement à la baisse des taux. Le
modèle n'est pas meilleur que la prévision naïve pour le PIB sur l'ensemble de la période. Les résultats à 2040-2047
reposent donc sur la cohérence théorique, pas sur une performance prédictive démontrée.
\item \textbf{La prime de risque} (0,03 point par point de dette) pèse autant que tous les comportements réunis sur la
dette de long terme ; elle n'était pas variée dans les configurations. Elle repose sur la littérature, pas sur une
estimation propre au projet.
\item \textbf{θ n'est pas identifié par les données françaises seules} : la sélection repose sur une durée d'écart de
production (dont l'intervalle est large) et sur un ancrage microéconomique étranger. Une estimation sur un panel de
pays serait nécessaire.
\item \textbf{La distribution dynamique est une approximation} : effets d'emploi en espérance, patrimoine non observé,
héritiers non identifiés, agents publics non identifiés (le point d'indice est réparti au prorata des salaires), une
seule évolution des salaires pour tous. Le résiduel de cohérence avec le compte macroéconomique reste d'environ un
point pour S3 (non attribué).
\item \textbf{Les effets d'offre de travail} des prestations (prime d'activité, RSA) sont dans le haut des évaluations
françaises ; ils pèsent sur les résultats de S1, S2 et S5.
\item \textbf{L'optimisation} repose sur un métamodèle additif non vérifié par le modèle complet pour la v0.8.
\item \textbf{Les comportements d'entreprises} (incidence de l'IS et des impôts de production, trésorerie,
délocalisation) restent paramétrés sur des études étrangères ou des hypothèses ; ils sont étiquetés comme tels.
\end{enumerate}

\section{Ce qu'un relecteur devrait vérifier en priorité}

La fermeture (équation et estimation de θ, ρ), le contrôle de cohérence micro-macro de M16 et ses clés de répartition,
la prime de risque, et la reproduction du choc d'inflation de 2022. Le code, les données téléchargées et tous les
résultats bruts sont dans le dépôt ; chaque table du rapport est produite par un script.
